\section{NeuS}
Tái tạo bề mặt 3D từ ảnh đa góc nhìn là một bài toán cốt lõi trong thị giác máy tính và đồ họa máy tính \cite{Wang2021NeuS}. Các phương pháp biểu diễn ngầm neural (neural implicit representations) gần đây cho thấy nhiều hứa hẹn nhờ khả năng tạo ra kết quả chất lượng cao và xử lý các trường hợp phức tạp như bề mặt không Lambertian hay cấu trúc mỏng \cite{Wang2021NeuS}. Tuy nhiên, các phương pháp hiện tại còn tồn tại những hạn chế nhất định.

\subsection{Hạn chế của các phương pháp trước NeuS}

Có hai hướng tiếp cận chính trước NeuS:

\begin{itemize}
    \item \textbf{Phương pháp dựa trên Surface Rendering (ví dụ: IDR \cite{Yariv2020IDR}, DVR \cite{Niemeyer2020DVR}):} Các phương pháp này thường biểu diễn bề mặt dưới dạng hàm chiếm dụng (occupancy) hoặc hàm khoảng cách có dấu (Signed Distance Function - SDF). Chúng sử dụng kỹ thuật surface rendering khả vi để huấn luyện mạng neural bằng cách so sánh ảnh render với ảnh đầu vào. Mặc dù IDR cho kết quả ấn tượng, nó gặp khó khăn với các cấu trúc phức tạp gây thay đổi độ sâu đột ngột. Nguyên nhân là do surface rendering chỉ xem xét một điểm giao duy nhất của tia nhìn với bề mặt, khiến gradient chỉ tồn tại cục bộ, dễ bị mắc kẹt ở cực tiểu địa phương. Thêm vào đó, các phương pháp này thường yêu cầu mặt nạ vật thể (foreground mask) để giám sát và hội tụ \cite{Wang2021NeuS}. % Source [2], [12-14], [16], [51]
    \item \textbf{Phương pháp dựa trên Volume Rendering (ví dụ: NeRF \cite{Mildenhall2020NeRF}):} Các phương pháp này biểu diễn cảnh dưới dạng trường bức xạ (radiance field) và mật độ thể tích (volume density). Chúng sử dụng volume rendering bằng cách lấy mẫu nhiều điểm dọc theo tia nhìn và tổng hợp màu sắc ($\alpha$-composition) để tạo ra màu pixel. Ưu điểm là khả năng xử lý thay đổi độ sâu đột ngột tốt hơn do gradient được lan truyền từ nhiều điểm. Tuy nhiên, NeRF được thiết kế chính cho tổng hợp ảnh nhìn mới (novel view synthesis), việc trích xuất bề mặt chất lượng cao từ trường mật độ học được là rất khó khăn do thiếu các ràng buộc hình học rõ ràng trên bề mặt. Bề mặt trích xuất từ NeRF thường bị nhiễu \cite{Wang2021NeuS}. % Source [3], [19-23], [52], [54]
\end{itemize}

\subsection{Phương pháp NeuS}

NeuS (Neural Surfaces) được đề xuất để khắc phục những hạn chế trên bằng cách kết hợp các ưu điểm của cả hai hướng tiếp cận: biểu diễn bề mặt chính xác bằng SDF và tối ưu mạnh mẽ bằng volume rendering \cite{Wang2021NeuS}. % Source [1], [29], [31], [55]

\subsubsection{Biểu diễn cảnh}
NeuS biểu diễn hình học của vật thể bằng một hàm SDF $f: \mathbb{R}^3 \to \mathbb{R}$ và biểu diễn màu sắc (appearance) bằng một hàm $c: \mathbb{R}^3 \times \mathbb{S}^2 \to \mathbb{R}^3$, cả hai đều được mã hóa bởi mạng MLP. Bề mặt $\mathcal{S}$ được định nghĩa là tập mức 0 của SDF: $\mathcal{S} = \{x \in \mathbb{R}^3 | f(x) = 0\}$ \cite{Wang2021NeuS}. % Source [5], [30], [60], [63], [64]

\subsubsection{Volume Rendering cho SDF}
Để áp dụng volume rendering cho việc huấn luyện mạng SDF, NeuS giới thiệu một hàm mật độ xác suất $\phi_s(f(x))$, gọi là \textbf{S-density}, được định nghĩa dựa trên giá trị SDF tại điểm $x$. Hàm $\phi_s$ thường được chọn là hàm mật độ logistic (đạo hàm của hàm sigmoid $\Phi_s(x)=(1+e^{-sx})^{-1}$), có dạng hình chuông tập trung tại $f(x)=0$. Tham số $1/s$ thể hiện độ lệch chuẩn và cũng là một tham số có thể học được \cite{Wang2021NeuS}. % Source [65-70]

Màu sắc của một pixel $C(o,v)$ được tính bằng tích phân màu $c(p(t), v)$ dọc theo tia nhìn $p(t) = o + tv$ với một hàm trọng số $w(t)$:
$$C(o,v) = \int_{0}^{+\infty} w(t) c(p(t), v) dt$$
\cite{Wang2021NeuS}. % Source [72], [73]

\subsubsection{Hàm Trọng số Volume Rendering Không Lệch (Unbiased)}
Việc xây dựng hàm trọng số $w(t)$ phù hợp là chìa khóa. NeuS chỉ ra rằng việc áp dụng công thức volume rendering tiêu chuẩn một cách "ngây thơ" sẽ dẫn đến sai lệch (bias) hình học. Để giải quyết vấn đề này, NeuS đề xuất một công thức volume rendering mới, đảm bảo tính \textbf{không lệch (unbiased)} (trọng số cực đại tại bề mặt SDF=0) và \textbf{nhận biết che khuất (occlusion-aware)} \cite{Wang2021NeuS}. % Source [6], [32], [74-80], [83-87]

NeuS định nghĩa một hàm "mật độ mờ đục" (opaque density) $\rho(t)$ mới từ S-density $\phi_s$ và CDF $\Phi_s$ của nó:
$$\rho(t) = \max\left(\frac{-\frac{d\Phi_s}{dt}(f(p(t)))}{\Phi_s(f(p(t)))}, 0\right)$$
và tính trọng số $w(t) = T(t)\rho(t)$. Hàm trọng số này được chứng minh là không lệch trong xấp xỉ bậc nhất \cite{Wang2021NeuS}. % Source [90-94], [101-103], [105-109], [111-114]

Trong quá trình rời rạc hóa, giá trị alpha $\alpha_i$ cho đoạn $[t_i, t_{i+1}]$ được tính bằng:
$$\alpha_i = \max\left(\frac{\Phi_s(f(p(t_i))) - \Phi_s(f(p(t_{i+1})))}{\Phi_s(f(p(t_i)))}, 0\right)$$
\cite{Wang2021NeuS}. % Source [115], [116], [263-266]

\subsubsection{Huấn luyện}
NeuS được huấn luyện bằng cách tối thiểu hóa hàm mất mát bao gồm $\mathcal{L}_{color}$ (L1 loss giữa màu render và ảnh gốc), $\mathcal{L}_{reg}$ (ràng buộc Eikonal: $\|\nabla f(x)\|_2 = 1$) và tùy chọn $\mathcal{L}_{mask}$ (nếu có mask) \cite{Wang2021NeuS}. % Source [117-119], [121-124]
NeuS cũng sử dụng chiến lược lấy mẫu phân cấp (hierarchical sampling) \cite{Wang2021NeuS}. % Source [124], [125]

\subsection{Kết quả và Ưu điểm}
Các thí nghiệm trên bộ dữ liệu DTU \cite{Jensen2014DTU} và BlendedMVS \cite{Yao2020BlendedMVS} cho thấy NeuS vượt trội hơn các phương pháp tiên tiến như IDR \cite{Yariv2020IDR}, NeRF \cite{Mildenhall2020NeRF}, COLMAP \cite{Schonberger16SFMRevisited} (hoặc \cite{Schoenberger2016Pixelwise}) và UNISURF \cite{Oechsle2021UNISURF} về chất lượng tái tạo bề mặt, đặc biệt đối với các đối tượng có cấu trúc phức tạp, tự che khuất và cấu trúc mỏng. Đáng chú ý, NeuS đạt được kết quả tốt ngay cả khi không cần sử dụng mask giám sát \cite{Wang2021NeuS}. % Source [7], [28], [35], [147], [149], [173], [178]

Tóm lại, NeuS là một phương pháp hiệu quả để tái tạo bề mặt 3D chất lượng cao từ ảnh đa góc nhìn, giải quyết được các hạn chế của các phương pháp trước đó bằng cách kết hợp biểu diễn SDF và một cơ chế volume rendering mới, không lệch.


\section{Instant-NGP}
Các phương pháp biểu diễn đồ họa dựa trên mạng neural (neural graphics primitives), thường sử dụng Multi-Layer Perceptrons (MLP) để tham số hóa các hàm biểu diễn hình dạng (ví dụ: SDF \cite{Park2019DeepSDF}) hoặc trường bức xạ (ví dụ: NeRF \cite{Mildenhall2020NeRF}), đã đạt được những thành tựu ấn tượng \cite{Muller2022InstantNGP}. Tuy nhiên, việc huấn luyện và tính toán các MLP lớn thường tốn kém tài nguyên và thời gian \cite{Muller2022InstantNGP_source_8}. Các kỹ thuật mã hóa đầu vào (input encoding) như mã hóa tần số (frequency encoding) \cite{Mildenhall2020NeRF, Vaswani2017Attention} giúp cải thiện chất lượng nhưng vẫn đòi hỏi MLP lớn và hội tụ chậm \cite{Muller2022InstantNGP_source_78}. Các phương pháp mã hóa tham số (parametric encoding) sử dụng cấu trúc dữ liệu phụ trợ như lưới dày (dense grid) hoặc cây (tree) \cite{Liu2020NSVF, Takikawa2021NGLOD} có thể tăng tốc huấn luyện bằng cách giảm kích thước MLP, nhưng lại gặp vấn đề về bộ nhớ lớn hoặc yêu cầu các thuật toán cập nhật cấu trúc phức tạp (pruning, merging), không hiệu quả trên GPU hiện đại \cite{Muller2022InstantNGP_source_26, Muller2022InstantNGP_source_28, Muller2022InstantNGP_source_55, Muller2022InstantNGP_source_74, Muller2022InstantNGP_source_81, Muller2022InstantNGP_source_92}.

Để giải quyết những thách thức này, Müller và cộng sự đã đề xuất Instant Neural Graphics Primitives (Instant-NGP), một phương pháp đột phá dựa trên kỹ thuật mã hóa đầu vào mới gọi là \textbf{Multiresolution Hash Encoding} (Mã hóa Băm Đa Phân giải) \cite{Muller2022InstantNGP}. Ý tưởng cốt lõi là thay thế MLP lớn bằng một MLP \textit{nhỏ} kết hợp với một cấu trúc dữ liệu hiệu quả chứa các tham số có thể huấn luyện được \cite{Muller2022InstantNGP_source_9}.

\subsection{Phương pháp Multiresolution Hash Encoding}

Mã hóa băm đa phân giải hoạt động như sau (xem Hình 3 trong \cite{Muller2022InstantNGP}):
\begin{enumerate}
    \item \textbf{Cấu trúc đa mức (Multiresolution Levels):} Sử dụng $L$ (thường là 16) mức phân giải khác nhau. Mỗi mức $l$ tương ứng với một lưới ảo (virtual grid) có độ phân giải $N_l$, tăng dần theo cấp số nhân từ $N_{min}$ đến $N_{max}$ \cite{Muller2022InstantNGP_source_112, Muller2022InstantNGP_source_113, Muller2022InstantNGP_source_114}.
    \item \textbf{Bảng băm đặc trưng (Feature Hash Tables):} Mỗi mức $l$ được liên kết với một bảng (mảng) chứa tối đa $T$ vector đặc trưng (feature vectors) $F$-chiều (thường $F=2$), gọi là $\theta_l$. $T$ là một siêu tham số quan trọng, cân bằng giữa chất lượng, bộ nhớ và hiệu năng \cite{Muller2022InstantNGP_source_105, Muller2022InstantNGP_source_112, Muller2022InstantNGP_source_125}.
    \item \textbf{Tra cứu và Băm (Lookup and Hashing):} Với một tọa độ đầu vào $x \in \mathbb{R}^d$, tại mỗi mức $l$:
        \begin{itemize}
            \item Xác định các đỉnh (corners) $2^d$ của ô lưới ảo (voxel) bao quanh $x$.
            \item Ánh xạ tọa độ nguyên của các đỉnh này tới các chỉ số (indices) trong bảng đặc trưng $\theta_l$:
                \begin{itemize}
                    \item Ở các mức phân giải thô (số đỉnh lưới $\le T$): Ánh xạ 1-1.
                    \item Ở các mức phân giải mịn (số đỉnh lưới $> T$): Sử dụng một hàm băm không gian (spatial hash function), ví dụ $h(\mathbf{x}) = (\bigoplus_{i=1}^d x_i \pi_i) \pmod T$, để lấy chỉ số \cite{Muller2022InstantNGP_source_104, Muller2022InstantNGP_source_117, Muller2022InstantNGP_source_118, Muller2022InstantNGP_source_121, Teschner2003SpatialHashing}.
                \end{itemize}
        \end{itemize}
    \item \textbf{Xử lý Xung đột Băm (Collision Handling):} Instant-NGP \textit{không} sử dụng các cơ chế xử lý xung đột băm tường minh. Thay vào đó, nó dựa vào hai yếu tố:
        \begin{itemize}
            \item Tối ưu hóa dựa trên gradient: Khi nhiều vị trí bị băm vào cùng một mục trong bảng, gradient của chúng sẽ được cộng dồn (hoặc trung bình). Gradient từ các mẫu "quan trọng" hơn (ví dụ: điểm trên bề mặt nhìn thấy trong NeRF) sẽ chiếm ưu thế và điều khiển giá trị của mục đó \cite{Muller2022InstantNGP_source_34, Muller2022InstantNGP_source_150, Muller2022InstantNGP_source_153}.
            \item Mạng MLP theo sau: Mạng neural nhỏ $m$ học cách phân giải các xung đột còn lại bằng cách kết hợp thông tin từ nhiều mức phân giải (các mức thô thường không có xung đột) \cite{Muller2022InstantNGP_source_99, Muller2022InstantNGP_source_119, Muller2022InstantNGP_source_142, Muller2022InstantNGP_source_149}.
        \end{itemize}
    \item \textbf{Nội suy và Kết hợp (Interpolation and Concatenation):} Các vector đặc trưng $F$-chiều được lấy ra từ bảng băm tại các đỉnh voxel được nội suy tuyến tính đa chiều (d-linear interpolation) dựa trên vị trí tương đối của $x$ trong voxel đó \cite{Muller2022InstantNGP_source_105, Muller2022InstantNGP_source_123}. Kết quả nội suy từ tất cả $L$ mức được nối (concatenate) lại với nhau (cùng các đầu vào phụ trợ $\xi$ như hướng nhìn nếu cần) để tạo thành vector đầu vào $y$ cho MLP $m(y; \Phi)$ \cite{Muller2022InstantNGP_source_106, Muller2022InstantNGP_source_124}.
\end{enumerate}

\subsection{Ưu điểm và Cài đặt}

Cách tiếp cận này mang lại nhiều lợi ích đáng kể:
\begin{itemize}
    \item \textbf{Tốc độ huấn luyện cực nhanh:} Cho phép huấn luyện các biểu diễn neural chất lượng cao chỉ trong vài giây đến vài phút trên một GPU duy nhất, nhanh hơn nhiều bậc độ lớn so với các phương pháp trước đó \cite{Muller2022InstantNGP_source_1, Muller2022InstantNGP_source_6, Muller2022InstantNGP_source_16, Muller2022InstantNGP_source_131, Muller2022InstantNGP_source_293, Muller2022InstantNGP_source_345}.
    \item \textbf{Hiệu năng tính toán cao:} Render thời gian thực (vài chục ms ở độ phân giải HD) \cite{Muller2022InstantNGP_source_16, Muller2022InstantNGP_source_255, Muller2022InstantNGP_source_272}.
    \item \textbf{Chất lượng cao:} Đạt được chất lượng tương đương hoặc tốt hơn các phương pháp SOTA như NeRF, mip-NeRF, NSVF, NGLOD trong nhiều tác vụ \cite{Muller2022InstantNGP_source_6, Muller2022InstantNGP_source_30, Muller2022InstantNGP_source_138, Muller2022InstantNGP_source_216, Muller2022InstantNGP_source_292}.
    \item \textbf{Tính linh hoạt và Đơn giản:} Mã hóa không phụ thuộc tác vụ, dễ cấu hình (chủ yếu qua $T$), không cần cơ chế cập nhật cấu trúc dữ liệu phức tạp \cite{Muller2022InstantNGP_source_7, Muller2022InstantNGP_source_29, Muller2022InstantNGP_source_36, Muller2022InstantNGP_source_94, Muller2022InstantNGP_source_195, Muller2022InstantNGP_source_343}. Tự động thích ứng với phân bố dữ liệu \cite{Muller2022InstantNGP_source_158, Muller2022InstantNGP_source_159}.
    \item \textbf{Hiệu quả trên GPU:} Cấu trúc bảng băm dễ song song hóa, tra cứu $O(1)$, tận dụng cache hiệu quả. Toàn bộ hệ thống được cài đặt bằng các CUDA kernel hợp nhất (fully-fused) trong framework \texttt{tiny-cuda-nn} \cite{Muller2021TinyCuDNN, Muller2022InstantNGP_source_15, Muller2022InstantNGP_source_37, Muller2022InstantNGP_source_38, Muller2022InstantNGP_source_39, Muller2022InstantNGP_source_100, Muller2022InstantNGP_source_167}.
\end{itemize}

Instant-NGP sử dụng MLP rất nhỏ (thường chỉ 2 lớp ẩn, 64 neuron) \cite{Muller2022InstantNGP_source_178}, lưu trữ đặc trưng dạng half-precision, tối ưu bằng Adam với $\epsilon$ nhỏ ($10^{-15}$) \cite{Kingma2014Adam, Muller2022InstantNGP_source_168, Muller2022InstantNGP_source_184}. Đối với NeRF, nhóm tác giả sử dụng kiến trúc gồm MLP mật độ và MLP màu, kết hợp occupancy grid để tăng tốc ray marching \cite{Muller2022InstantNGP_source_262, Muller2022InstantNGP_source_270, Muller2022InstantNGP_source_490}.

\subsection{Ứng dụng và Hạn chế}

Instant-NGP đã được chứng minh hiệu quả trên nhiều tác vụ như biểu diễn ảnh gigapixel, học hàm SDF, Neural Radiance Caching \cite{Muller2021NRC}, và đặc biệt là NeRF \cite{Muller2022InstantNGP_source_1}. Nó cho phép huấn luyện NeRF chất lượng cao nhanh hơn 20-60 lần so với frequency encoding trên cùng implementation tối ưu \cite{Muller2022InstantNGP_source_281, Muller2022InstantNGP_source_299}.

Một hạn chế nhỏ là hiện tượng nhiễu hạt (microstructure/graininess) có thể xuất hiện do xung đột băm, đặc biệt rõ trên bề mặt SDF \cite{Muller2022InstantNGP_source_228, Muller2022InstantNGP_source_245, Muller2022InstantNGP_source_328}.

Nhìn chung, Instant-NGP với mã hóa băm đa phân giải là một bước tiến quan trọng, cung cấp một giải pháp thay thế hiệu quả, nhanh chóng và linh hoạt cho các cấu trúc dữ liệu và MLP lớn trong lĩnh vực neural graphics primitives.
\section{NeuS2}
Chắc chắn rồi, dựa trên bài báo bạn cung cấp về NeuS2, tôi đã tổng hợp thông tin tập trung vào phương pháp, ưu nhược điểm và chuẩn bị một section LaTeX cùng file BibTeX tương ứng.

1. Mã LaTeX cho Section trong Luận văn:

Đoạn mã

\section{NeuS2: Học Nhanh Bề Mặt Ngầm Neural để Tái tạo Đa Góc nhìn}

Các phương pháp tái tạo bề mặt 3D dựa trên biểu diễn ngầm neural, như NeuS \cite{Wang2021NeuS}, đã chứng minh khả năng tạo ra kết quả chất lượng rất cao cho cảnh tĩnh \cite{Wang2023NeuS2_source_1}. Tuy nhiên, một hạn chế lớn của NeuS là thời gian huấn luyện cực kỳ dài (khoảng 8 giờ), khiến nó gần như không thể áp dụng cho các cảnh động với hàng nghìn khung hình \cite{Wang2023NeuS2_source_2, Wang2023NeuS2_source_17}. Mặt khác, các phương pháp tăng tốc như Instant-NGP \cite{Muller2022InstantNGP} tuy nhanh nhưng chất lượng hình học trích xuất từ trường mật độ (density field) lại bị nhiễu và thiếu các ràng buộc bề mặt \cite{Wang2023NeuS2_source_19, Wang2023NeuS2_source_94}. Việc kết hợp trực tiếp NeuS và Instant-NGP gặp thách thức lớn trong việc tính toán hiệu quả đạo hàm bậc hai (cần thiết cho ràng buộc Eikonal của NeuS) khi sử dụng các MLP được tối ưu hóa cao bằng CUDA \cite{Wang2023NeuS2_source_23, Wang2023NeuS2_source_24, Wang2023NeuS2_source_95, Wang2023NeuS2_source_96, Wang2023NeuS2_source_97, Wang2023NeuS2_source_98}.

NeuS2 \cite{Wang2023NeuS2} được đề xuất để giải quyết những vấn đề này, nhằm mục tiêu học nhanh các bề mặt ngầm neural chất lượng cao từ ảnh đa góc nhìn cho cả cảnh tĩnh và động.

\subsection{Phương pháp Mô hình NeuS2}

NeuS2 kế thừa ý tưởng biểu diễn bề mặt bằng hàm khoảng cách có dấu (SDF) và cơ chế volume rendering không lệch (unbiased) từ NeuS \cite{Wang2021NeuS}, đồng thời tích hợp các kỹ thuật tăng tốc từ Instant-NGP \cite{Muller2022InstantNGP} và đề xuất các cải tiến mới.

\subsubsection{Kiến trúc cho Cảnh tĩnh}
\begin{itemize}
    \item \textbf{Biểu diễn kết hợp SDF và Hash Encoding:} Thay vì dùng MLP lớn, NeuS2 sử dụng một MLP \textit{nhỏ} để học hàm SDF $f$. Đầu vào của MLP này là sự kết hợp (concatenate) của tọa độ điểm 3D $x$ và vector đặc trưng $h_\Omega(x)$ được truy vấn từ cấu trúc mã hóa băm đa phân giải (Multi-resolution Hash Encoding) với các mục trong bảng băm $\Omega$ có thể học được \cite{Wang2023NeuS2_source_4, Wang2023NeuS2_source_22, Wang2023NeuS2_source_86, Wang2023NeuS2_source_110, Wang2023NeuS2_source_112}. Mạng SDF này dự đoán giá trị SDF $d$ và một vector đặc trưng hình học $g$ \cite{Wang2023NeuS2_source_112}.
    \item \textbf{Mạng Màu:} Một MLP nhỏ khác dự đoán màu sắc $c$, nhận đầu vào là vị trí $x$, pháp tuyến $n = \nabla_x d$, hướng nhìn $v$, giá trị SDF $d$, và đặc trưng hình học $g$ \cite{Wang2023NeuS2_source_114, Wang2023NeuS2_source_115}.
    \item \textbf{Volume Rendering:} Sử dụng công thức volume rendering không lệch của NeuS \cite{Wang2021NeuS} để tính màu pixel từ SDF và màu dự đoán dọc theo tia nhìn \cite{Wang2023NeuS2_source_116, Wang2023NeuS2_source_405, Wang2023NeuS2_source_409}. Quá trình ray marching được tăng tốc bằng occupancy grid tương tự Instant-NGP \cite{Wang2023NeuS2_source_117, Wang2023NeuS2_source_411, Wang2023NeuS2_source_412}.
    \item \textbf{Tính toán Đạo hàm Bậc hai Hiệu quả:} Đây là đóng góp kỹ thuật quan trọng của NeuS2 để xử lý ràng buộc Eikonal $\mathcal{L}_{eikonal} = \sum (\|\nabla_x d\| - 1)^2$ một cách hiệu quả \cite{Wang2023NeuS2_source_119, Wang2023NeuS2_source_120}. Thay vì dùng finite difference (dễ gây mất ổn định \cite{Wang2023NeuS2_source_99, Wang2023NeuS2_source_178}) hoặc autograd của các framework (chậm \cite{Wang2023NeuS2_source_124}), NeuS2 đưa ra công thức giải tích \textit{chính xác} và rút gọn cho đạo hàm bậc hai của hàm SDF $d$ đối với các tham số của mạng SDF và bảng băm, đặc biệt tận dụng tính chất của MLP dùng ReLU ($\frac{\partial^2 y}{\partial x^2} = 0$) \cite{Wang2023NeuS2_source_4, Wang2023NeuS2_source_25, Wang2023NeuS2_source_32, Wang2023NeuS2_source_100, Wang2023NeuS2_source_125, Wang2023NeuS2_source_129, Wang2023NeuS2_source_130}. Điều này cho phép cài đặt hiệu quả bằng CUDA kernel, tăng tốc đáng kể quá trình backpropagation \cite{Wang2023NeuS2_source_4, Wang2023NeuS2_source_132}.
    \item \textbf{Huấn luyện Tiến bộ (Progressive Training):} Để cải thiện sự ổn định và chất lượng khi dùng hash encoding, NeuS2 áp dụng chiến lược huấn luyện tăng dần độ phân giải. Các mức trong bảng băm được kích hoạt tuần tự từ thô đến mịn trong quá trình huấn luyện thông qua một hệ số trọng số phụ thuộc vào tiến trình huấn luyện \cite{Wang2023NeuS2_source_5, Wang2023NeuS2_source_26, Wang2023NeuS2_source_33, Wang2023NeuS2_source_109, Wang2023NeuS2_source_135, Wang2023NeuS2_source_136, Wang2023NeuS2_source_137}.
\end{itemize}

\subsubsection{Kiến trúc cho Cảnh động}
NeuS2 mở rộng sang tái tạo cảnh động bằng hai thành phần chính:
\begin{itemize}
    \item \textbf{Huấn luyện Tăng cường (Incremental Training):} Thay vì huấn luyện mỗi frame độc lập, NeuS2 tận dụng tính liên tục thời gian. Frame đầu tiên được huấn luyện từ đầu, các frame tiếp theo được khởi tạo từ các tham số (bảng băm, trọng số MLP) đã học của frame trước đó và chỉ cần fine-tune trong thời gian ngắn (~20 giây/frame). Điều này giúp tăng tốc đáng kể quá trình học cho toàn bộ chuỗi video \cite{Wang2023NeuS2_source_6, Wang2023NeuS2_source_28, Wang2023NeuS2_source_103, Wang2023NeuS2_source_140, Wang2023NeuS2_source_142, Wang2023NeuS2_source_143}.
    \item \textbf{Dự đoán Biến đổi Toàn cục (Global Transformation Prediction - GTP):} Khi có chuyển động lớn giữa các frame, việc fine-tune trực tiếp có thể gặp vấn đề (mắc kẹt ở cực tiểu địa phương). NeuS2 đề xuất dự đoán một phép biến đổi cứng (xoay $R_i$, tịnh tiến $T_i$) giữa frame hiện tại $i$ và frame trước đó $i-1$. Các điểm 3D trong frame $i$ được ánh xạ về không gian gốc (canonical space - thường là frame 0) thông qua chuỗi biến đổi tích lũy trước khi được đưa vào mạng SDF/màu. GTP giúp ổn định quá trình huấn luyện tăng cường và cho phép mô hình hóa các chuỗi dài với chuyển động lớn trong một vùng không gian giới hạn, tiết kiệm bộ nhớ \cite{Wang2023NeuS2_source_6, Wang2023NeuS2_source_30, Wang2023NeuS2_source_34, Wang2023NeuS2_source_90, Wang2023NeuS2_source_104, Wang2023NeuS2_source_145, Wang2023NeuS2_source_146, Wang2023NeuS2_source_148, Wang2023NeuS2_source_150, Wang2023NeuS2_source_151, Wang2023NeuS2_source_152}. GTP được tối ưu đồng thời trong quá trình fine-tune \cite{Wang2023NeuS2_source_163}.
\end{itemize}

\subsection{Ưu điểm và Nhược điểm}

\subsubsection{Ưu điểm}
\begin{itemize}
    \item \textbf{Tốc độ Vượt trội:} NeuS2 nhanh hơn NeuS khoảng hai bậc độ lớn (ví dụ: 5 phút so với 8 giờ cho cảnh tĩnh trên DTU) \cite{Wang2023NeuS2_source_3, Wang2023NeuS2_source_10, Wang2023NeuS2_source_50, Wang2023NeuS2_source_157}. Tái tạo cảnh động chỉ mất khoảng 20 giây cho mỗi frame (sau frame đầu tiên) \cite{Wang2023NeuS2_source_11, Wang2023NeuS2_source_21, Wang2023NeuS2_source_197, Wang2023NeuS2_source_208}.
    \item \textbf{Chất lượng Tái tạo Cao:} Đạt được chất lượng hình học và bề ngoài tương đương hoặc tốt hơn các phương pháp SOTA như NeuS, Instant-NGP, Instant-NSR, Voxurf, D-NeRF, TiNeuVox, đặc biệt là hình học chi tiết và ít nhiễu hơn Instant-NGP \cite{Wang2023NeuS2_source_3, Wang2023NeuS2_source_7, Wang2023NeuS2_source_9, Wang2023NeuS2_source_20, Wang2023NeuS2_source_53, Wang2023NeuS2_source_173, Wang2023NeuS2_source_181, Wang2023NeuS2_source_190, Wang2023NeuS2_source_194, Wang2023NeuS2_source_199, Wang2023NeuS2_source_209}.
    \item \textbf{Khả năng Xử lý Cảnh động Phức tạp:} Kiến trúc huấn luyện tăng cường kết hợp GTP cho phép xử lý hiệu quả các chuỗi video dài (hàng nghìn frame) với chuyển động lớn và biến dạng phức tạp, vượt qua giới hạn của nhiều phương pháp trước đó \cite{Wang2023NeuS2_source_6, Wang2023NeuS2_source_27, Wang2023NeuS2_source_34, Wang2023NeuS2_source_60, Wang2023NeuS2_source_62, Wang2023NeuS2_source_104, Wang2023NeuS2_source_152, Wang2023NeuS2_source_183, Wang2023NeuS2_source_195, Wang2023NeuS2_source_204, Wang2023NeuS2_source_205}.
    \item \textbf{Huấn luyện Ổn định:} Các kỹ thuật như tính đạo hàm bậc hai chính xác và huấn luyện tiến bộ giúp quá trình học ổn định hơn so với các phương pháp xấp xỉ hoặc huấn luyện trực tiếp trên độ phân giải cao \cite{Wang2023NeuS2_source_5, Wang2023NeuS2_source_100, Wang2023NeuS2_source_133, Wang2023NeuS2_source_178}.
\end{itemize}

\subsubsection{Nhược điểm và Hạn chế}
\begin{itemize}
    \item \textbf{Thiếu Đối ứng Dày đặc (Dense Correspondences):} Đối với cảnh động, NeuS2 tái tạo hình học cho từng frame một cách hiệu quả nhưng không cung cấp trực tiếp thông tin đối ứng dày đặc giữa các điểm trên bề mặt qua thời gian \cite{Wang2023NeuS2_source_221}. Việc này có thể cần các bước xử lý hậu kỳ như tracking lưới \cite{Wang2023NeuS2_source_222}.
    \item \textbf{Chi phí Lưu trữ cho Cảnh động:} Do lưu trữ tham số (chủ yếu là bảng băm, khoảng 25MB/frame) cho mỗi frame, việc lưu trữ các chuỗi video rất dài có thể tốn kém \cite{Wang2023NeuS2_source_232}. Cần các kỹ thuật nén trong tương lai \cite{Wang2023NeuS2_source_233}.
    \item \textbf{Phụ thuộc vào ảnh đa góc nhìn và hiệu chỉnh máy ảnh:} Giống như các phương pháp dựa trên NeRF/NeuS khác, NeuS2 yêu cầu đầu vào là ảnh/video từ nhiều góc nhìn với thông số máy ảnh đã biết.
    \item \textbf{Sử dụng Mặt nạ (Masks):} Các thí nghiệm trong bài báo (ví dụ trên DTU) sử dụng mặt nạ đối tượng \cite{Wang2023NeuS2_source_167, Wang2023NeuS2_source_349}. Mặc dù không được nêu rõ là bắt buộc, mặt nạ có thể cần thiết hoặc giúp cải thiện chất lượng/tốc độ hội tụ.
\end{itemize}

Tóm lại, NeuS2 đại diện cho một bước tiến đáng kể trong việc tái tạo bề mặt neural, cân bằng xuất sắc giữa tốc độ, chất lượng và khả năng ứng dụng cho cả cảnh tĩnh và động phức tạp.

\section{Geo-Neus}
Các phương pháp học bề mặt ngầm neural thông qua volume rendering, điển hình là NeuS \cite{Wang2021NeuS} và VolSDF \cite{Yariv2021VolSDF}, đã cho thấy tiềm năng lớn trong việc tái tạo bề mặt 3D từ ảnh đa góc nhìn \cite{Fu2022GeoNeus_source_1}. Tuy nhiên, một thách thức cốt lõi vẫn tồn tại: các phương pháp này chủ yếu dựa vào việc tối ưu hàm mất mát màu sắc (rendering loss) và thiếu các ràng buộc hình học đa góc nhìn (multi-view geometry constraints) một cách tường minh \cite{Fu2022GeoNeus_source_2, Fu2022GeoNeus_source_15, Fu2022GeoNeus_source_33}. Điều này dẫn đến sự không nhất quán về mặt hình học trong bề mặt tái tạo, đặc biệt là ở các vùng ít texture hoặc cấu trúc phức tạp \cite{Fu2022GeoNeus_source_2, Fu2022GeoNeus_source_16}.

Fu và cộng sự \cite{Fu2022GeoNeus} chỉ ra rằng có một "khoảng cách" hay "sai lệch" (gap/bias) cố hữu giữa việc tối ưu tích phân màu sắc trong volume rendering và việc mô hình hóa bề mặt chính xác tại tập mức 0 (zero-level set) của hàm khoảng cách có dấu (SDF) \cite{Fu2022GeoNeus_source_4, Fu2022GeoNeus_source_18, Fu2022GeoNeus_source_19, Fu2022GeoNeus_source_68, Fu2022GeoNeus_source_81, Fu2022GeoNeus_source_87}. Tích phân volume rendering tối ưu trên nhiều điểm dọc tia nhìn, trong khi bề mặt thực sự chỉ được xác định tại một điểm duy nhất (SDF=0). Sự sai lệch này khiến việc chỉ dựa vào mất mát màu sắc để giám sát mạng SDF trở nên không chính xác và khó khăn trong việc tối ưu hình học thực sự \cite{Fu2022GeoNeus_source_24, Fu2022GeoNeus_source_88}.

Để giải quyết vấn đề này, Geo-Neus \cite{Fu2022GeoNeus} được đề xuất với mục tiêu học bề mặt ngầm neural nhất quán về mặt hình học bằng cách đưa vào các ràng buộc hình học đa góc nhìn một cách tường minh.

\subsection{Phương pháp và Mô hình Geo-Neus}

Geo-Neus xây dựng dựa trên nền tảng của NeuS, sử dụng mạng MLP để biểu diễn SDF ($F_\Theta$) và trường bức xạ ($G_\Phi$), cùng với kỹ thuật volume rendering không lệch \cite{Fu2022GeoNeus_source_66, Fu2022GeoNeus_source_78}. Điểm cốt lõi của Geo-Neus là bổ sung hai loại ràng buộc hình học tường minh trực tiếp lên mạng SDF, thay vì chỉ dựa vào giám sát gián tiếp qua mất mát màu sắc \cite{Fu2022GeoNeus_source_3, Fu2022GeoNeus_source_29, Fu2022GeoNeus_source_37, Fu2022GeoNeus_source_90, Fu2022GeoNeus_source_91} (Xem Hình 2 trong \cite{Fu2022GeoNeus}).

\begin{enumerate}
    \item \textbf{Ràng buộc SDF từ Điểm Thưa (Sparse 3D Points - $\mathcal{L}_{SDF}$):}
        \begin{itemize}
            \item \textbf{Nguồn dữ liệu:} Tận dụng các điểm 3D thưa $P$ được tạo ra như một sản phẩm phụ "miễn phí" từ quá trình Structure from Motion (SfM) \cite{Schonberger16SFMRevisited, Snavely2006PhotoTourism} khi ước lượng tham số máy ảnh \cite{Fu2022GeoNeus_source_30, Fu2022GeoNeus_source_94, Fu2022GeoNeus_source_95, Fu2022GeoNeus_source_96}.
            \item \textbf{Giả định:} Các điểm SfM này được coi là nằm trên bề mặt thực của vật thể, do đó giá trị SDF tại các điểm này lý tưởng là bằng 0 ($sdf(p_i) \approx 0$) \cite{Fu2022GeoNeus_source_97, Fu2022GeoNeus_source_98}.
            \item \textbf{Xử lý che khuất:} Với mỗi ảnh đầu vào $I_i$, chỉ sử dụng các điểm SfM $P_i$ nhìn thấy được từ góc nhìn đó để giám sát (dựa trên liên kết với các điểm đặc trưng 2D $X_i$) \cite{Fu2022GeoNeus_source_100, Fu2022GeoNeus_source_101, Fu2022GeoNeus_source_103}.
            \item \textbf{Hàm mất mát:} $\mathcal{L}_{SDF} = \frac{1}{N_i} \sum_{p \in P_i} \| s\hat{d}f(p) \|_1$, với $s\hat{d}f(p)$ là giá trị SDF dự đoán bởi mạng $F_\Theta$. Mất mát này trực tiếp ép buộc mạng SDF dự đoán giá trị 0 tại các điểm đã biết trên bề mặt, đặc biệt hiệu quả ở những vùng có texture phong phú nơi SfM hoạt động tốt \cite{Fu2022GeoNeus_source_104, Fu2022GeoNeus_source_105, Fu2022GeoNeus_source_169}.
        \end{itemize}

    \item \textbf{Ràng buộc Nhất quán Quang trắc Đa góc nhìn (Multi-view Photometric Consistency - $\mathcal{L}_{photo}$):}
        \begin{itemize}
            \item \textbf{Mục tiêu:} Đảm bảo tính nhất quán hình học của bề mặt được biểu diễn bởi SDF=0 trên các góc nhìn khác nhau, đặc biệt hữu ích cho các vùng bề mặt lớn, trơn, ít texture nơi điểm SfM thưa thớt \cite{Fu2022GeoNeus_source_30, Fu2022GeoNeus_source_107, Fu2022GeoNeus_source_108, Fu2022GeoNeus_source_169}.
            \item \textbf{Xác định Điểm Bề mặt trên Tia nhìn:} Thay vì dùng tích phân độ sâu như trong volume rendering (vốn bị bias \cite{Fu2022GeoNeus_source_174, Fu2022GeoNeus_source_178}), Geo-Neus trực tiếp xác định vị trí $t^*$ trên tia nhìn $p(t) = o+tv$ nơi hàm SDF dự đoán $s\hat{d}f(t)$ đổi dấu và bằng 0 (sử dụng nội suy tuyến tính giữa các điểm lấy mẫu rời rạc $t_i, t_{i+1}$ có $s\hat{d}f(t_i) \cdot s\hat{d}f(t_{i+1}) < 0$) \cite{Fu2022GeoNeus_source_31, Fu2022GeoNeus_source_109, Fu2022GeoNeus_source_110, Fu2022GeoNeus_source_114, Fu2022GeoNeus_source_117, Fu2022GeoNeus_source_118}. Điểm giao đầu tiên $t^*$ được chọn để đảm bảo tính nhìn thấy \cite{Fu2022GeoNeus_source_119, Fu2022GeoNeus_source_120}.
            \item \textbf{Áp dụng Nguyên lý MVS:} Với điểm bề mặt $p = o+t^*v$ và pháp tuyến $n = \nabla s\hat{d}f(p)$ tại đó, sử dụng phép biến đổi homography phẳng (plane-induced homography) $H$ để ánh xạ (warp) một vùng ảnh nhỏ (patch) xung quanh điểm chiếu của $p$ trên ảnh tham chiếu $I_r$ sang các ảnh nguồn $I_s$ \cite{Fu2022GeoNeus_source_122, Fu2022GeoNeus_source_125, Fu2022GeoNeus_source_126, Fu2022GeoNeus_source_127, Hartley2004MVG}.
            \item \textbf{Đo lường sự nhất quán:} Tính toán sự tương đồng quang trắc giữa patch trên ảnh tham chiếu và các patch được ánh xạ trên ảnh nguồn bằng cách sử dụng Normalized Cross-Correlation (NCC) trên ảnh xám \cite{Fu2022GeoNeus_source_128, Fu2022GeoNeus_source_130, Furukawa2010MVS, Galliani2015Gipuma}.
            \item \textbf{Hàm mất mát:} $\mathcal{L}_{photo}$ được định nghĩa dựa trên $1 - NCC$ của các patch tương ứng, thường lấy trung bình của một số (ví dụ: 4) giá trị NCC tốt nhất để xử lý che khuất \cite{Fu2022GeoNeus_source_132}. Mất mát này khuyến khích hình học (vị trí $p$ và pháp tuyến $n$) phải nhất quán giữa các góc nhìn khác nhau \cite{Fu2022GeoNeus_source_133}.
        \end{itemize}

    \item \textbf{Hàm Mất mát Tổng thể:} Kết hợp các thành phần:
        $$\mathcal{L} = \mathcal{L}_{color} + \alpha\mathcal{L}_{reg} + \beta\mathcal{L}_{SDF} + \gamma\mathcal{L}_{photo}$$
        trong đó $\mathcal{L}_{color}$ là mất mát màu sắc (L1 loss), $\mathcal{L}_{reg}$ là ràng buộc Eikonal ($\|\nabla s\hat{d}f\| \approx 1$) \cite{Gropp2020IGR}, $\mathcal{L}_{SDF}$ và $\mathcal{L}_{photo}$ là các mất mát hình học mới. Các hệ số $\alpha, \beta, \gamma$ cân bằng các thành phần \cite{Fu2022GeoNeus_source_134}.
\end{enumerate}

Kiến trúc mạng MLP cho SDF và màu sắc tương tự như trong NeuS \cite{Wang2021NeuS} và IDR \cite{Yariv2020IDR}, sử dụng positional encoding \cite{Mildenhall2020NeRF} và khởi tạo hình học \cite{Atzmon2020SAL} \cite{Fu2022GeoNeus_source_145, Fu2022GeoNeus_source_146, Fu2022GeoNeus_source_147}.

\subsection{Phân tích Ưu nhược điểm}

\subsubsection{Ưu điểm}
\begin{itemize}
    \item \textbf{Chất lượng Hình học Vượt trội:} Nhờ các ràng buộc hình học đa góc nhìn tường minh ($\mathcal{L}_{SDF}$, $\mathcal{L}_{photo}$), Geo-Neus tạo ra các bề mặt có độ chính xác và nhất quán hình học cao hơn đáng kể so với các phương pháp chỉ dựa vào giám sát màu sắc như NeuS, VolSDF \cite{Fu2022GeoNeus_source_2, Fu2022GeoNeus_source_3, Fu2022GeoNeus_source_7, Fu2022GeoNeus_source_32, Fu2022GeoNeus_source_34, Fu2022GeoNeus_source_133, Fu2022GeoNeus_source_154, Fu2022GeoNeus_source_155, Fu2022GeoNeus_source_187}.
    \item \textbf{Giảm Sai lệch (Bias) của Volume Rendering:} Việc tối ưu trực tiếp trên bề mặt SDF=0 thông qua các ràng buộc hình học giúp giảm thiểu tác động của sai lệch vốn có trong phương pháp tích phân màu của volume rendering \cite{Fu2022GeoNeus_source_4, Fu2022GeoNeus_source_6, Fu2022GeoNeus_source_31, Fu2022GeoNeus_source_38, Fu2022GeoNeus_source_87, Fu2022GeoNeus_source_178}.
    \item \textbf{Tái tạo Tốt Cấu trúc Phức tạp và Vùng Trơn:} Sự kết hợp của $\mathcal{L}_{SDF}$ (tập trung vào vùng có texture, cấu trúc mỏng) và $\mathcal{L}_{photo}$ (tập trung vào vùng trơn, ít texture) giúp Geo-Neus tái tạo tốt cả hai loại bề mặt này \cite{Fu2022GeoNeus_source_7, Fu2022GeoNeus_source_34, Fu2022GeoNeus_source_40, Fu2022GeoNeus_source_156, Fu2022GeoNeus_source_169, Fu2022GeoNeus_source_170, Fu2022GeoNeus_source_186}.
    \item \textbf{Hội tụ Nhanh hơn:} Giám sát hình học tường minh giúp mạng hội tụ nhanh hơn và ổn định hơn so với NeuS \cite{Fu2022GeoNeus_source_105, Fu2022GeoNeus_source_180, Fu2022GeoNeus_source_182}.
\end{itemize}

\subsubsection{Nhược điểm và Hạn chế}
\begin{itemize}
    \item \textbf{Chi phí Tính toán và Thời gian Huấn luyện:} Mặc dù hội tụ nhanh hơn về số vòng lặp, việc tính toán thêm các ràng buộc hình học (đặc biệt là $\mathcal{L}_{photo}$ với warping và NCC) làm tăng chi phí tính toán trên mỗi vòng lặp. Thời gian huấn luyện tổng thể vẫn còn đáng kể (đề cập khoảng 16 giờ), tương đương NeuS và chậm hơn nhiều so với các phương pháp tăng tốc như Instant-NGP hay NeuS2 \cite{Fu2022GeoNeus_source_149, Fu2022GeoNeus_source_188}.
    \item \textbf{Phụ thuộc vào SfM:} Chất lượng và độ phủ của các điểm 3D thưa từ SfM có thể ảnh hưởng đến hiệu quả của $\mathcal{L}_{SDF}$. Các vùng hoàn toàn không có texture sẽ không nhận được giám sát từ loại ràng buộc này.
    \item \textbf{Độ phức tạp:} Mô hình và quy trình huấn luyện phức tạp hơn do có thêm các thành phần xử lý hình học (SfM, MVS).
    \item \textbf{Không đề cập đến Cảnh động:} Bài báo chỉ tập trung vào tái tạo cảnh tĩnh.
\end{itemize}

Tóm lại, Geo-Neus là một cải tiến quan trọng cho các phương pháp học bề mặt ngầm neural bằng cách tích hợp các ràng buộc hình học đa góc nhìn tường minh, giúp khắc phục sai lệch của volume rendering và cải thiện đáng kể chất lượng, độ nhất quán của bề mặt tái tạo, dù vẫn còn hạn chế về tốc độ tính toán.

\section{Neuralangelo}
Các phương pháp tái tạo bề mặt 3D dựa trên mạng nơ-ron, đặc biệt là những phương pháp sử dụng biểu diễn hình học ẩn (implicit geometric representations) kết hợp với kết xuất thể tích khả vi (differentiable volume rendering), đã cho thấy tiềm năng to lớn trong việc khôi phục các bề mặt dày đặc và hoàn chỉnh từ các tập ảnh đa góc nhìn \cite{Wang21NeuS, Yariv21VolSDF, Yariv20IDR}. Tuy nhiên, một thách thức lớn còn tồn tại là khả năng tái tạo các chi tiết hình học tinh vi (fine-grained details) và mở rộng quy mô cho các cảnh thực tế lớn và phức tạp \cite{Li23Neuralangelo}. Các phương pháp trước đó thường dựa vào mạng MLP tiêu chuẩn hoặc mã hóa tần số (frequency/positional encoding) \cite{Mildenhall2020NeRF, Tancik20Fourier}, vốn gặp khó khăn trong việc biểu diễn hiệu quả các tín hiệu tần số cao trên quy mô lớn \cite{Muller22InstantNGP}.

Để giải quyết những hạn chế này, Li và cộng sự đã đề xuất \textbf{Neuralangelo} \cite{Li23Neuralangelo}, một phương pháp mới nhằm đạt được khả năng tái tạo bề mặt 3D với độ trung thực cao từ ảnh RGB thông thường. Neuralangelo kết hợp sức mạnh biểu diễn của **lưới băm đa độ phân giải (multi-resolution hash grids)**, được giới thiệu trong Instant NGP \cite{Muller22InstantNGP}, với kỹ thuật kết xuất bề mặt nơ-ron dựa trên **Hàm khoảng cách có dấu (Signed Distance Function - SDF)**. Hai thành phần kỹ thuật then chốt được giới thiệu để khai thác tối đa tiềm năng của lưới băm trong bài toán tái tạo bề mặt là: (1) việc sử dụng \textbf{gradient số (numerical gradients)} để tính toán các đạo hàm bậc cao, giúp ổn định quá trình tối ưu hóa và tăng cường tính trơn tru; và (2) một chiến lược **tối ưu hóa từ thô đến tinh (coarse-to-fine optimization)** trên các cấp độ phân giải của lưới băm, cho phép kiểm soát và tái tạo hiệu quả các chi tiết ở các quy mô khác nhau.

\subsection{Nền tảng: Biểu diễn ẩn và Lưới băm}
\label{subsec:neuralangelo_background}

Neuralangelo xây dựng dựa trên các kỹ thuật biểu diễn bề mặt 3D ẩn và kết xuất thể tích khả vi. Cảnh 3D được biểu diễn bởi một hàm SDF $f(\vect{x}): \mathbb{R}^3 \rightarrow \mathbb{R}$, trong đó bề mặt $\mathcal{S}$ là tập hợp các điểm có giá trị SDF bằng không: $\mathcal{S} = \{\vect{x} \in \mathbb{R}^3 | f(\vect{x}) = 0\}$. Hàm SDF này được xấp xỉ bởi một mạng nơ-ron $f_\Theta$. Ngoài ra, một mạng nơ-ron khác $c_\Theta(\vect{x}, \vect{v}, \vect{n}, \vect{z})$ được sử dụng để dự đoán màu sắc (RGB) phụ thuộc vào vị trí $\vect{x}$, hướng nhìn $\vect{v}$, pháp tuyến bề mặt $\vect{n} = \nabla f(\vect{x})$, và một vector đặc trưng tiềm ẩn $\vect{z}$ từ mạng SDF \cite{Yariv20IDR}.

Để kết xuất màu sắc $\hat{c}$ của một tia nhìn từ các hàm SDF và màu này, các phương pháp như NeuS \cite{Wang21NeuS} và VolSDF \cite{Yariv21VolSDF} đề xuất các cách chuyển đổi giá trị SDF thành mật độ thể tích $\sigma$ hoặc độ chắn sáng $\alpha$ để sử dụng trong công thức kết xuất thể tích (ví dụ, công thức (3) trong bài báo Neuralangelo \cite{Li23Neuralangelo}, dựa trên \cite{Wang21NeuS}). Màu sắc pixel cuối cùng được tính bằng cách tích hợp thông tin màu và mật độ/opacity dọc theo tia (công thức (1) trong bài báo \cite{Li23Neuralangelo}). Quá trình tối ưu hóa thường bao gồm việc giảm thiểu sự khác biệt giữa màu pixel kết xuất $\hat{c}$ và màu pixel thực tế $c$ (ví dụ, sử dụng $\mathcal{L}_1$ loss, công thức (2) \cite{Li23Neuralangelo}), kết hợp với các ràng buộc hình học như Eikonal loss $\mathcal{L}_{eik}$ để khuyến khích $||\nabla f(\vect{x})||_2 = 1$.

Điểm khác biệt cốt lõi của Neuralangelo nằm ở cách biểu diễn đầu vào cho mạng nơ-ron SDF và màu sắc. Thay vì sử dụng trực tiếp tọa độ $\vect{x}$ hoặc mã hóa tần số Fourier \cite{Tancik20Fourier}, Neuralangelo sử dụng **mã hóa lưới băm đa độ phân giải (multi-resolution hash encoding)** $\gamma(\vect{x})$ từ Instant NGP \cite{Muller22InstantNGP}. Kỹ thuật này sử dụng $L$ lưới có độ phân giải khác nhau $V_l$, từ thô đến tinh. Với một điểm $\vect{x}$, đặc trưng tại mỗi cấp độ phân giải $l$ được lấy bằng cách nội suy tuyến tính (trilinear interpolation) các vector đặc trưng lưu trữ tại các đỉnh của ô lưới (grid cell) chứa điểm đó. Các vector đặc trưng này được lấy từ một bảng băm (hash table) có kích thước giới hạn, giúp kiểm soát bộ nhớ hiệu quả. Các đặc trưng từ $L$ cấp độ phân giải sau đó được nối lại $\gamma(\vect{x}) = (\gamma_1(\vect{x}), ..., \gamma_L(\vect{x}))$ và đưa vào một mạng MLP nông (shallow MLP) để dự đoán SDF và đặc trưng màu $\vect{z}$ [71-77]. Cách biểu diễn này cho phép mã hóa các chi tiết ở nhiều tần số khác nhau một cách hiệu quả về bộ nhớ và tính toán [81].

\subsection{Phương pháp Neuralangelo}
\label{subsec:neuralangelo_method}

Neuralangelo kế thừa kiến trúc biểu diễn SDF dựa trên lưới băm đa độ phân giải và kết xuất thể tích tương tự NeuS/VolSDF, nhưng giới thiệu hai cải tiến quan trọng để khắc phục các vấn đề gặp phải khi áp dụng trực tiếp lưới băm vào bài toán tái tạo bề mặt SDF.

\subsubsection{Sử dụng Gradient số cho Đạo hàm bậc cao}
\label{subsubsec:neuralangelo_numerical_grad}

Việc áp đặt các ràng buộc hình học như Eikonal loss ($\mathcal{L}_{eik} = \mathbb{E}[(||\nabla f(\vect{x})||_2 - 1)^2]$, công thức (5) \cite{Li23Neuralangelo}) đòi hỏi tính toán gradient của SDF, $\nabla f(\vect{x})$. Các phương pháp trước đây thường sử dụng gradient giải tích (analytical gradients) thông qua backpropagation \cite{Wang21NeuS, Yariv20IDR, Yariv21VolSDF, Yang21Geometry}. Tuy nhiên, các tác giả Neuralangelo chỉ ra rằng gradient giải tích của đặc trưng lưới băm $\gamma(\vect{x})$ đối với vị trí $\vect{x}$ (khi dùng nội suy tuyến tính) có tính cục bộ và không liên tục tại biên giới các ô lưới (công thức (7) \cite{Li23Neuralangelo}) [91-94]. Điều này có nghĩa là gradient của Eikonal loss (vốn yêu cầu đạo hàm bậc hai của $f$ đối với trọng số mạng) chỉ lan truyền đến các mục trong bảng băm (hash entries) tương ứng với các đỉnh của ô lưới chứa điểm $\vect{x}$ đang xét [95]. Điều này cản trở việc học các bề mặt trơn tru và nhất quán trải rộng qua nhiều ô lưới, vì thông tin gradient không được chia sẻ hiệu quả giữa các vùng lân cận.

Để giải quyết vấn đề này, Neuralangelo đề xuất tính toán $\nabla f(\vect{x})$ (và các đạo hàm bậc cao hơn nếu cần, như cho độ cong) bằng phương pháp **gradient số (numerical gradients)** thông qua sai phân hữu hạn (finite differences) [84]:
\begin{equation} \label{eq:numerical_gradient}
\nabla_x f(\vect{x}) \approx \frac{f(\gamma(\vect{x} + \epsilon_x)) - f(\gamma(\vect{x} - \epsilon_x))}{2\epsilon}
\end{equation}
trong đó $\epsilon_x = [\epsilon, 0, 0]^T$ và các thành phần khác được tính tương tự [102]. Bước nhảy $\epsilon$ đóng vai trò quan trọng. Khi tính toán gradient số, mạng SDF $f$ được truy vấn tại các điểm $\vect{x} \pm \epsilon_k$. Nếu $\epsilon$ đủ lớn (ví dụ, lớn hơn kích thước ô lưới), các truy vấn này sẽ lấy đặc trưng từ nhiều ô lưới khác nhau. Do đó, khi lan truyền ngược gradient của Eikonal loss (hoặc curvature loss), các cập nhật sẽ ảnh hưởng đến các mục băm trong một vùng rộng hơn thay vì chỉ các đỉnh của ô lưới cục bộ [96]. Về bản chất, gradient số với bước nhảy $\epsilon$ phù hợp hoạt động như một phép toán làm trơn (smoothing operation) trên trường gradient giải tích cục bộ và không liên tục của lưới băm, giúp tối ưu hóa các bề mặt nhất quán hơn [97, 100].

\subsubsection{Tối ưu hóa Từ Thô đến Tinh}
\label{subsubsec:neuralangelo_coarse_to_fine}

Để tái tạo hiệu quả các chi tiết bề mặt ở các cấp độ khác nhau, Neuralangelo áp dụng một chiến lược tối ưu hóa **từ thô đến tinh (coarse-to-fine)**, được điều khiển bởi cả hai yếu tố: bước nhảy $\epsilon$ của gradient số và việc kích hoạt các cấp độ phân giải của lưới băm [106, 107].

\begin{itemize}
    \item \textbf{Giảm dần bước nhảy $\epsilon$ (Annealing $\epsilon$):} Như đã đề cập, $\epsilon$ kiểm soát mức độ "nhìn xa" của gradient số và mức độ làm trơn bề mặt. Neuralangelo bắt đầu quá trình tối ưu hóa với giá trị $\epsilon$ lớn, tương ứng với kích thước ô lưới ở cấp độ phân giải thô nhất. Sau đó, $\epsilon$ được giảm dần theo hàm mũ trong quá trình huấn luyện, tiến tới các giá trị nhỏ hơn tương ứng với các cấp độ lưới mịn hơn [111]. Khi $\epsilon$ lớn, Eikonal loss và curvature loss sẽ khuyến khích các bề mặt trơn tru trên quy mô lớn. Khi $\epsilon$ nhỏ dần, các ràng buộc này sẽ tập trung vào các chi tiết cục bộ hơn, cho phép mạng học các cấu trúc tinh vi [108-110].
    \item \textbf{Kích hoạt lưới băm phân cấp (Progressive Hash Grid Activation):} Thay vì kích hoạt tất cả $L$ cấp độ phân giải của lưới băm ngay từ đầu, Neuralangelo chỉ bắt đầu với một tập hợp các cấp độ thô nhất. Khi quá trình tối ưu hóa diễn ra và bước nhảy $\epsilon$ giảm xuống tương ứng với kích thước ô lưới của một cấp độ phân giải $l$ nào đó, cấp độ đó mới được "kích hoạt" và bắt đầu đóng góp vào đặc trưng $\gamma(\vect{x})$ [113]. Các cấp độ mịn hơn sẽ được kích hoạt tuần tự sau đó. Chiến lược này giúp mạng tập trung học các cấu trúc tổng thể ở giai đoạn đầu với các lưới thô, sau đó mới học các chi tiết tinh vi hơn với các lưới mịn khi chúng được kích hoạt. Điều này tránh việc các lưới mịn phải "học lại" hoặc bị ảnh hưởng tiêu cực bởi các cập nhật ở giai đoạn tối ưu hóa cấu trúc thô, giúp bảo toàn chi tiết tốt hơn [112, 115, 116].
\end{itemize}
Ngoài ra, kỹ thuật suy giảm trọng số (weight decay) cũng được áp dụng để tránh việc một vài cấp độ phân giải nào đó chiếm ưu thế quá lớn trong kết quả cuối cùng [117].

\subsubsection{Hàm mất mát tổng thể}
\label{subsubsec:neuralangelo_loss}

Bên cạnh Eikonal loss $\mathcal{L}_{eik}$ (công thức (5)) và color loss $\mathcal{L}_{RGB}$ (công thức (2)), Neuralangelo còn sử dụng một thành phần **ràng buộc độ cong (Curvature Regularization)** $\mathcal{L}_{curv}$ để khuyến khích bề mặt trơn tru hơn nữa [118]. Độ cong trung bình (mean curvature) được xấp xỉ bằng toán tử Laplace số $\nabla^2 f(\vect{x})$, cũng được tính bằng sai phân hữu hạn từ các điểm đã lấy mẫu để tính gradient số [119, 121]. Hàm mất mát độ cong được định nghĩa là:
\begin{equation} \label{eq:curvature_loss}
\mathcal{L}_{curv} = \frac{1}{N} \sum_{i=1}^{N} |\nabla^2 f(\vect{x}_i)|.
\end{equation}
Hàm mất mát tổng thể của Neuralangelo là tổng có trọng số của ba thành phần:
\begin{equation} \label{eq:total_loss}
\mathcal{L} = \mathcal{L}_{RGB} + w_{eik}\mathcal{L}_{eik} + w_{curv}\mathcal{L}_{curv},
\end{equation}
với $w_{eik}$ và $w_{curv}$ là các trọng số điều chỉnh tầm quan trọng tương đối của các ràng buộc hình học [121]. Một chiến lược "warm-up" cho $w_{curv}$ cũng được áp dụng để tránh làm phẳng các cấu trúc lõm ở giai đoạn đầu tối ưu hóa [173-176]. Toàn bộ các tham số mạng (MLPs và các mục trong bảng băm) được tối ưu hóa đồng thời end-to-end [122].

\subsection{Phân tích Ưu điểm và Nhược điểm}
\label{subsec:neuralangelo_pros_cons}

Neuralangelo đã chứng minh được hiệu quả vượt trội so với các phương pháp tái tạo bề mặt nơ-ron trước đó, đặc biệt trên các bộ dữ liệu quy mô lớn và phức tạp như Tanks and Temples \cite{Knapitsch17Tanks}.

\textbf{Ưu điểm:}
\begin{itemize}
    \item \textbf{Độ trung thực cao (High Fidelity):} Nhờ sử dụng lưới băm đa độ phân giải và chiến lược tối ưu hóa từ thô đến tinh, Neuralangelo có khả năng tái tạo các chi tiết hình học bề mặt rất tinh vi, vượt xa các phương pháp dựa trên MLP tiêu chuẩn hoặc mã hóa tần số Fourier [143, 149].
    \item \textbf{Khả năng mở rộng (Scalability):} Tận dụng hiệu quả về bộ nhớ và tính toán của lưới băm \cite{Muller22InstantNGP}, Neuralangelo có thể xử lý các cảnh quy mô lớn (ví dụ: các tòa nhà, khu vực ngoài trời) mà các phương pháp trước đó gặp khó khăn [156, 202].
    \item \textbf{Không yêu cầu dữ liệu phụ trợ:} Phương pháp hoạt động hiệu quả chỉ với ảnh RGB và pose máy ảnh đầu vào, không cần thông tin về độ sâu (depth) hay mặt nạ phân đoạn (segmentation masks) như một số phương pháp khác \cite{Li23Neuralangelo_Abstract, Fu22GeoNeuS, Yu22MonoSDF} [7, 53].
    \item \textbf{Chất lượng bề mặt tốt:} Việc sử dụng gradient số và ràng buộc độ cong giúp tạo ra các bề mặt trơn tru và ít nhiễu hơn so với việc trích xuất bề mặt trực tiếp từ trường mật độ của NeRF hoặc các phương pháp SDF sử dụng gradient giải tích trên lưới băm [137, 141, 180].
\end{itemize}

\textbf{Nhược điểm và Hạn chế:}
\begin{itemize}
    \item \textbf{Thời gian huấn luyện:} Mặc dù hiệu quả hơn MLP sâu, việc huấn luyện Neuralangelo cho mỗi cảnh vẫn đòi hỏi thời gian đáng kể (hàng giờ hoặc hơn tùy thuộc vào kích thước cảnh và phần cứng) [186].
    \item \textbf{Nhạy cảm với siêu tham số:} Hiệu suất của phương pháp phụ thuộc vào việc lựa chọn cẩn thận các siêu tham số, đặc biệt là lịch trình giảm $\epsilon$ và kích hoạt các cấp độ lưới băm, cũng như các trọng số $w_{eik}, w_{curv}$.
    \item \textbf{Vấn đề với bề mặt phản xạ mạnh:} Bài báo thừa nhận rằng, giống như Instant NGP, Neuralangelo có thể gặp khó khăn hơn so với các phương pháp dựa trên mã hóa tần số Fourier (như NeRF gốc) khi xử lý các bề mặt có độ phản xạ rất cao (highly reflective surfaces), đây là một hạn chế kế thừa từ bản chất của mã hóa băm [235, 236].
    \item \textbf{Yêu cầu Pose máy ảnh chính xác:} Giống như hầu hết các phương pháp dựa trên NeRF và SfM, Neuralangelo yêu cầu đầu vào là các pose máy ảnh đã được ước lượng tương đối chính xác (ví dụ, từ COLMAP \cite{Schonberger16SFMRevisited}) [200]. Lỗi lớn trong pose máy ảnh đầu vào sẽ ảnh hưởng tiêu cực đến chất lượng tái tạo.
    \item \textbf{Chiến lược lấy mẫu điểm/pixel:} Bài báo gốc đề cập rằng việc lấy mẫu pixel ngẫu nhiên hiện tại chưa tối ưu và việc khám phá các chiến lược lấy mẫu hiệu quả hơn là một hướng nghiên cứu trong tương lai để tăng tốc độ huấn luyện [185-187].
\end{itemize}

Tóm lại, Neuralangelo đại diện cho một bước tiến quan trọng trong lĩnh vực tái tạo bề mặt 3D từ ảnh đa góc nhìn, đặc biệt là đối với các cảnh phức tạp và quy mô lớn, bằng cách kết hợp hiệu quả biểu diễn lưới băm đa độ phân giải với các kỹ thuật tối ưu hóa và xử lý gradient phù hợp cho bài toán tái tạo bề mặt SDF.

